\documentclass[10pt,a4paper]{amsart}
\usepackage{amsmath,amssymb,amsthm,graphicx}
\usepackage{verbatim}
\newenvironment{reference}{\comment}{\endcomment}
%\newenvironment{reference}{}{}
\newenvironment{slogan}{\comment}{\endcomment}
\newenvironment{history}{\comment}{\endcomment}
% For commutative diagrams we use Xy-pic
\usepackage[all]{xy}

% We use 2cell for 2-commutative diagrams.
\xyoption{2cell}
% Theorem environments.
%
\theoremstyle{plain}
\newtheorem{theorem}[subsection]{Theorem}
\newtheorem{proposition}[subsection]{Proposition}
\newtheorem{lemma}[subsection]{Lemma}

\theoremstyle{definition}
\newtheorem{definition}[subsection]{Definition}
\newtheorem{example}[subsection]{Example}
\newtheorem{exercise}[subsection]{Exercise}
\newtheorem{situation}[subsection]{Situation}

\theoremstyle{remark}
\newtheorem{remark}[subsection]{Remark}
\newtheorem{remarks}[subsection]{Remarks}

\numberwithin{equation}{subsection}

% Macros
%
\def\lim{\mathop{\mathrm{lim}}\nolimits}
\def\colim{\mathop{\mathrm{colim}}\nolimits}
\def\Spec{\mathop{\mathrm{Spec}}}
\def\Hom{\mathop{\mathrm{Hom}}\nolimits}
\def\Ext{\mathop{\mathrm{Ext}}\nolimits}
\def\SheafHom{\mathop{\mathcal{H}\!\mathit{om}}\nolimits}
\def\SheafExt{\mathop{\mathcal{E}\!\mathit{xt}}\nolimits}
\def\Sch{\mathit{Sch}}
\def\Mor{\mathop{\mathrm{Mor}}\nolimits}
\def\Ob{\mathop{\mathrm{Ob}}\nolimits}
\def\Sh{\mathop{\mathit{Sh}}\nolimits}
\def\NL{\mathop{N\!L}\nolimits}
\def\CH{\mathop{\mathrm{CH}}\nolimits}
\def\proetale{{pro\text{-}\acute{e}tale}}
\def\etale{{\acute{e}tale}}
\def\QCoh{\mathit{QCoh}}
\def\Ker{\mathop{\mathrm{Ker}}}
\def\Im{\mathop{\mathrm{Im}}}
\def\Coker{\mathop{\mathrm{Coker}}}
\def\Coim{\mathop{\mathrm{Coim}}}

% Boxtimes
%
\DeclareMathSymbol{\boxtimes}{\mathbin}{AMSa}{"02}

%
% Macros for moduli stacks/spaces
%
\def\QCohstack{\mathcal{QC}\!\mathit{oh}}
\def\Cohstack{\mathcal{C}\!\mathit{oh}}
\def\Spacesstack{\mathcal{S}\!\mathit{paces}}
\def\Quotfunctor{\mathrm{Quot}}
\def\Hilbfunctor{\mathrm{Hilb}}
\def\Curvesstack{\mathcal{C}\!\mathit{urves}}
\def\Polarizedstack{\mathcal{P}\!\mathit{olarized}}
\def\Complexesstack{\mathcal{C}\!\mathit{omplexes}}
% \Pic is the operator that assigns to X its picard group, usage \Pic(X)
% \Picardstack_{X/B} denotes the Picard stack of X over B
% \Picardfunctor_{X/B} denotes the Picard functor of X over B
\def\Pic{\mathop{\mathrm{Pic}}\nolimits}
\def\Picardstack{\mathcal{P}\!\mathit{ic}}
\def\Picardfunctor{\mathrm{Pic}}
\def\Deformationcategory{\mathcal{D}\!\mathit{ef}}
\usepackage{cleveref}
\crefname{originalreference}{original reference}{original references}
\Crefname{originalreference}{Original reference}{Original references}
\makeatletter
\newlabel{algebraic-section-stack-to-presentation}{{algebraic-section-stack-to-presentation}{source}{}{sourcepointer-1}{}}
\newlabel{algebraic-section-stack-to-presentation@cref}{{[originalreference][0][]algebraic-section-stack-to-presentation}{[1][0][]source}{}{}{}}
\makeatother
\crefname{originalreference}{original reference}{original references}
\Crefname{originalreference}{Original reference}{Original references}
\makeatletter
\newlabel{criteria-theorem-flat-groupoid-gives-algebraic-stack}{{criteria-theorem-flat-groupoid-gives-algebraic-stack}{source}{}{sourcepointer-2}{}}
\newlabel{criteria-theorem-flat-groupoid-gives-algebraic-stack@cref}{{[originalreference][0][]criteria-theorem-flat-groupoid-gives-algebraic-stack}{[1][0][]source}{}{}{}}
\makeatother
\crefname{originalreference}{original reference}{original references}
\Crefname{originalreference}{Original reference}{Original references}
\makeatletter
\newlabel{spaces-groupoids-lemma-quotient-stack-restrict}{{spaces-groupoids-lemma-quotient-stack-restrict}{source}{}{sourcepointer-3}{}}
\newlabel{spaces-groupoids-lemma-quotient-stack-restrict@cref}{{[originalreference][0][]spaces-groupoids-lemma-quotient-stack-restrict}{[1][0][]source}{}{}{}}
\makeatother
\crefname{originalreference}{original reference}{original references}
\Crefname{originalreference}{Original reference}{Original references}
\makeatletter
\newlabel{spaces-morphisms-lemma-closed-immersion-finite-presentation}{{spaces-morphisms-lemma-closed-immersion-finite-presentation}{source}{}{sourcepointer-4}{}}
\newlabel{spaces-morphisms-lemma-closed-immersion-finite-presentation@cref}{{[originalreference][0][]spaces-morphisms-lemma-closed-immersion-finite-presentation}{[1][0][]source}{}{}{}}
\makeatother
\crefname{originalreference}{original reference}{original references}
\Crefname{originalreference}{Original reference}{Original references}
\makeatletter
\newlabel{spaces-morphisms-lemma-composition-finite-presentation}{{spaces-morphisms-lemma-composition-finite-presentation}{source}{}{sourcepointer-5}{}}
\newlabel{spaces-morphisms-lemma-composition-finite-presentation@cref}{{[originalreference][0][]spaces-morphisms-lemma-composition-finite-presentation}{[1][0][]source}{}{}{}}
\makeatother
\crefname{originalreference}{original reference}{original references}
\Crefname{originalreference}{Original reference}{Original references}
\makeatletter
\newlabel{spaces-morphisms-lemma-locally-finite-type-locally-noetherian}{{spaces-morphisms-lemma-locally-finite-type-locally-noetherian}{source}{}{sourcepointer-6}{}}
\newlabel{spaces-morphisms-lemma-locally-finite-type-locally-noetherian@cref}{{[originalreference][0][]spaces-morphisms-lemma-locally-finite-type-locally-noetherian}{[1][0][]source}{}{}{}}
\makeatother
\crefname{originalreference}{original reference}{original references}
\Crefname{originalreference}{Original reference}{Original references}
\makeatletter
\newlabel{stacks-properties-lemma-monomorphism}{{stacks-properties-lemma-monomorphism}{source}{}{sourcepointer-7}{}}
\newlabel{stacks-properties-lemma-monomorphism@cref}{{[originalreference][0][]stacks-properties-lemma-monomorphism}{[1][0][]source}{}{}{}}
\makeatother
\crefname{originalreference}{original reference}{original references}
\Crefname{originalreference}{Original reference}{Original references}
\makeatletter
\newlabel{stacks-properties-lemma-open-substacks}{{stacks-properties-lemma-open-substacks}{source}{}{sourcepointer-8}{}}
\newlabel{stacks-properties-lemma-open-substacks@cref}{{[originalreference][0][]stacks-properties-lemma-open-substacks}{[1][0][]source}{}{}{}}
\makeatother
\crefname{originalreference}{original reference}{original references}
\Crefname{originalreference}{Original reference}{Original references}
\makeatletter
\newlabel{stacks-properties-section-residual-gerbe}{{stacks-properties-section-residual-gerbe}{source}{}{sourcepointer-9}{}}
\newlabel{stacks-properties-section-residual-gerbe@cref}{{[originalreference][0][]stacks-properties-section-residual-gerbe}{[1][0][]source}{}{}{}}
\makeatother
\title{Every point of a locally Noetherian algebraic stack has a residual gerbe}
\author{The Stacks Project Authors}
\date{stacks-project:stacks-morphisms.tex:local-snapshot-20261012; GNU FDL 1.2 or later}
\begin{document}
\maketitle
\noindent\textit{Editorial note.} All delivered mathematical text is the author's own, copied from the pinned Stacks Project chapter source sources/stacks-algebraic-stacks-morphisms/pinned-original-source.tex; no statement or proof step was rewritten, reordered or removed.\par\medskip
\noindent\textit{Editorial note.} Preamble provenance: the Stacks Project keeps one single shared preamble for every chapter; this card takes it directly from the staged sources/stacks-algebraic-stacks-morphisms/preamble.tex as three verbatim spans (lines 8-12, 14-18 and 154-217), so no derived file was written and nothing was added or redefined. The wrapper class and the base packages are supplied by the assembler.\par\medskip
\noindent\textit{Editorial note.} Bibliography provenance: neither the counted statement nor its proof contains a citation (this lane rescans the whole span before writing the card), so the card carries no bibliography and no entry was reconstructed.\par\medskip
\noindent\textit{Editorial note.} Selection provenance: this card was opened under the file-based control contract (CONTROL.md), with the mechanical pre-card gate run on its statement and proof interval before the card was written.\par\medskip
\noindent\textit{Allowed standard prerequisites (editorial).}\begin{itemize}
\item The presentation of an algebraic stack by a smooth groupoid in algebraic spaces with an affine Noetherian scheme of representatives. Invoked early in the counted proof (source lines 7172-7180) as pointers to results of the chapter on algebraic stacks. Provides the presentation in which the point and its residue field are realised.
\item The choice of a point of a Noetherian topological space minimal in the fibre over a given point, and the description of such a point as a limit of affine opens of the closure. Used inside the counted proof (source lines 7183-7206). Makes the inverse image of the closure meet the fibre exactly in the chosen point, which forces the factorisation used in the proof.
\item The restriction of a groupoid along a morphism and the resulting monomorphism of quotient stacks, and the identification of the residual gerbe through the uniqueness statement. Invoked at the end of the counted proof (source lines 7215-7229) as cross-chapter pointers. Shows that the quotient of the restricted groupoid is the residual gerbe.
\end{itemize}\medskip
\ifdefined\hypertarget\hypertarget{sourcepointer-1}{}\fi
\noindent\textit{Source reference.} Original reference algebraic-section-stack-to-presentation: sources/stacks-algebraic-stacks-morphisms/pinned-original-source.tex, source lines 6365--6365. Navigation pointer only: the counted text refers to the Stacks item algebraic-section-stack-to-presentation, which belongs to another chapter of the project and is not part of this card. The pointer carries the target's own label as its number because this chapter's staged snapshot does not contain the other chapter's tag table; its content is not reproduced here.\par\medskip
\ifdefined\hypertarget\hypertarget{sourcepointer-2}{}\fi
\noindent\textit{Source reference.} Original reference criteria-theorem-flat-groupoid-gives-algebraic-stack: sources/stacks-algebraic-stacks-morphisms/pinned-original-source.tex, source lines 3850--3850. Navigation pointer only: the counted text refers to the Stacks item criteria-theorem-flat-groupoid-gives-algebraic-stack, which belongs to another chapter of the project and is not part of this card. The pointer carries the target's own label as its number because this chapter's staged snapshot does not contain the other chapter's tag table; its content is not reproduced here.\par\medskip
\ifdefined\hypertarget\hypertarget{sourcepointer-3}{}\fi
\noindent\textit{Source reference.} Original reference spaces-groupoids-lemma-quotient-stack-restrict: sources/stacks-algebraic-stacks-morphisms/pinned-original-source.tex, source lines 7138--7138. Navigation pointer only: the counted text refers to the Stacks item spaces-groupoids-lemma-quotient-stack-restrict, which belongs to another chapter of the project and is not part of this card. The pointer carries the target's own label as its number because this chapter's staged snapshot does not contain the other chapter's tag table; its content is not reproduced here.\par\medskip
\ifdefined\hypertarget\hypertarget{sourcepointer-4}{}\fi
\noindent\textit{Source reference.} Original reference spaces-morphisms-lemma-closed-immersion-finite-presentation: sources/stacks-algebraic-stacks-morphisms/pinned-original-source.tex, source lines 7213--7213. Navigation pointer only: the counted text refers to the Stacks item spaces-morphisms-lemma-closed-immersion-finite-presentation, which belongs to another chapter of the project and is not part of this card. The pointer carries the target's own label as its number because this chapter's staged snapshot does not contain the other chapter's tag table; its content is not reproduced here.\par\medskip
\ifdefined\hypertarget\hypertarget{sourcepointer-5}{}\fi
\noindent\textit{Source reference.} Original reference spaces-morphisms-lemma-composition-finite-presentation: sources/stacks-algebraic-stacks-morphisms/pinned-original-source.tex, source lines 5624--5624. Navigation pointer only: the counted text refers to the Stacks item spaces-morphisms-lemma-composition-finite-presentation, which belongs to another chapter of the project and is not part of this card. The pointer carries the target's own label as its number because this chapter's staged snapshot does not contain the other chapter's tag table; its content is not reproduced here.\par\medskip
\ifdefined\hypertarget\hypertarget{sourcepointer-6}{}\fi
\noindent\textit{Source reference.} Original reference spaces-morphisms-lemma-locally-finite-type-locally-noetherian: sources/stacks-algebraic-stacks-morphisms/pinned-original-source.tex, source lines 7212--7212. Navigation pointer only: the counted text refers to the Stacks item spaces-morphisms-lemma-locally-finite-type-locally-noetherian, which belongs to another chapter of the project and is not part of this card. The pointer carries the target's own label as its number because this chapter's staged snapshot does not contain the other chapter's tag table; its content is not reproduced here.\par\medskip
\ifdefined\hypertarget\hypertarget{sourcepointer-7}{}\fi
\noindent\textit{Source reference.} Original reference stacks-properties-lemma-monomorphism: sources/stacks-algebraic-stacks-morphisms/pinned-original-source.tex, source lines 1167--1167. Navigation pointer only: the counted text refers to the Stacks item stacks-properties-lemma-monomorphism, which belongs to another chapter of the project and is not part of this card. The pointer carries the target's own label as its number because this chapter's staged snapshot does not contain the other chapter's tag table; its content is not reproduced here.\par\medskip
\ifdefined\hypertarget\hypertarget{sourcepointer-8}{}\fi
\noindent\textit{Source reference.} Original reference stacks-properties-lemma-open-substacks: sources/stacks-algebraic-stacks-morphisms/pinned-original-source.tex, source lines 7176--7176. Navigation pointer only: the counted text refers to the Stacks item stacks-properties-lemma-open-substacks, which belongs to another chapter of the project and is not part of this card. The pointer carries the target's own label as its number because this chapter's staged snapshot does not contain the other chapter's tag table; its content is not reproduced here.\par\medskip
\ifdefined\hypertarget\hypertarget{sourcepointer-9}{}\fi
\noindent\textit{Source reference.} Original reference stacks-properties-section-residual-gerbe: sources/stacks-algebraic-stacks-morphisms/pinned-original-source.tex, source lines 7229--7229. Navigation pointer only: the counted text refers to the Stacks item stacks-properties-section-residual-gerbe, which belongs to another chapter of the project and is not part of this card. The pointer carries the target's own label as its number because this chapter's staged snapshot does not contain the other chapter's tag table; its content is not reproduced here.\par\medskip
\begin{lemma}
\label{lemma-every-point-residual-gerbe-locally-Noetherian}
Let $\mathcal{X}$ be a locally Noetherian algebraic stack.
Then the residual gerbe of $\mathcal{X}$ at $x$ exists for
every $x \in |\mathcal{X}|$.
\end{lemma}

\begin{proof}
Choose an affine scheme $U$ and a smooth morphism $U \to \mathcal{X}$
such that $x$ is in the image of the open continuous map
$|U| \to |\mathcal{X}|$. We may and do replace $\mathcal{X}$
with the open substack corresponding to the image of $|U| \to |\mathcal{X}|$,
see Properties of Stacks, Lemma \ref{stacks-properties-lemma-open-substacks}.
Thus we may assume $\mathcal{X} = [U/R]$ for a smooth
groupoid $(U, R, s, t, c)$ in algebraic spaces
where $U$ is a Noetherian affine scheme, see
Algebraic Stacks, Section \ref{algebraic-section-stack-to-presentation}.

\medskip\noindent
Let $E \subset |U|$ be the inverse image of $\{x\} \subset |\mathcal{X}|$.
Of course $E \not = \emptyset$.
Since $|U|$ is a Noetherian topological space, we can choose an element
$u \in E$ such that $\overline{\{u\}} \cap E = \{u\}$.
As usual, we think of $u = \Spec(\kappa(u))$ as the spectrum of
its residue field. Let us write
$$
F = u \times_{U, t} R = u \times_\mathcal{X} U
\quad\text{and}\quad
R' = (u \times u) \times_{(U \times U), (t, s)} R = u \times_\mathcal{X} u
$$
Furthermore, denote $Z = \overline{\{u\}} \subset U$ with the reduced
induced scheme structure. Denote $p : F \to U$ the morphism induced
by the second projection (using $s : R \to U$ in the first fibre product
description of $F$). Then $E$ is the set theoretic image of $p$.
The morphism $R' \to F$ is a monomorphism which factors through the
inverse image $p^{-1}(Z)$ of $Z$. This inverse image
$p^{-1}(Z) \subset F$ is a closed subscheme and
the restriction $p|_{p^{-1}(Z)} : p^{-1}(Z) \to Z$ has image
set theoretically contained in $\{u\} \subset Z$ by our careful
choice of $u \in E$ above. Since $u = \lim W$ where the limit
is over the nonempty affine open subschemes of the irreducible reduced scheme
$Z$, we conclude that the morphism $p|_{p^{-1}(Z)} : p^{-1}(Z) \to Z$
factors through the morphism $u \to Z$. Clearly this implies that
$R' = p^{-1}(Z)$. In particular the morphism $t' : R' \to u$ is
locally of finite presentation as the composition of the closed
immersion $p^{-1}(Z) \to F$ of locally Noetherian algebraic spaces
with the smooth morphism $\text{pr}_1 : F \to u$; use
Morphisms of Spaces, Lemmas
\ref{spaces-morphisms-lemma-locally-finite-type-locally-noetherian},
\ref{spaces-morphisms-lemma-closed-immersion-finite-presentation}, and
\ref{spaces-morphisms-lemma-composition-finite-presentation}.
Hence the restriction $(u, R', s', t', c')$ of $(U, R, s, t, c)$
by $u \to U$ is a groupoid in algebraic spaces where $s'$ and $t'$
are flat and locally of finite presentation. Therefore
$\mathcal{Z} = [u/R']$ is an algebraic stack by
Criteria for Representability, Theorem
\ref{criteria-theorem-flat-groupoid-gives-algebraic-stack}.
As $R'$ is the restriction of $R$ by $u \to U$ we see
$\mathcal{Z} \to \mathcal{X}$ is a monomorphism by
Groupoids in Spaces, Lemma
\ref{spaces-groupoids-lemma-quotient-stack-restrict}
and
Properties of Stacks, Lemma \ref{stacks-properties-lemma-monomorphism}.
Then $\mathcal{Z}$ is (isomorphic to) the residual gerbe by
the material in Properties of Stacks, Section
\ref{stacks-properties-section-residual-gerbe}.
\end{proof}

\end{document}
